\section{Related Work}
\label{sec:related-work}

Our work sits at the intersection of three research traditions. This section
places it relative to each of them and states, plainly, what we do and do not
adopt from the mathematical vocabulary we borrow.

\paragraph{Knowledge graphs and semantic stores.}
The RDF triple model~\cite{manola2004rdf} and its query language
SPARQL~\cite{sparql2013} form the backbone of the modern semantic web. Production
triple stores---Apache Jena, Virtuoso, GraphDB, Blazegraph---have evolved highly
tuned index structures (SPO, POS, OSP, and permutations), custom storage
formats, and cost-based query optimisers for the reified triple representation.
Named graphs~\cite{carroll2005named} extend the triple with a fourth graph
coordinate and are the standard mechanism for scoping. Property
graphs~\cite{angles2017property}, as implemented in Neo4j and Amazon Neptune,
decorate edges with attribute maps, which handles some higher-arity cases
gracefully but reintroduces the reification problem for facts of arity above
two or three (Section~\ref{sec:related-work:lpg} returns to this comparison
in more depth below).

The reification cost of representing $n$-ary facts in triples is well
documented~\cite{angles2008survey,noy2006nary}. It is $O(k)$ triples and
$O(k)$ joins per fact of arity $k$. Our KG baseline implements the standard
reified approach with SPO/POS/OSP indexes and dictionary encoding. We do
not claim this baseline is state of the art; it is a controlled reference
point.

\paragraph{Statement-level annotation and $n$-ary mitigation.}
The reification tax has motivated a family of mitigations inside the
triple model itself. RDF-star~\cite{hartig2017rdfstar,w3c2021rdfstar}
allows a triple to appear as the subject or object of another triple, so
statement-level metadata (provenance, confidence, temporal scope) no
longer needs a minted event node; the singleton-property
encoding~\cite{nguyen2014singleton} achieves a similar effect by minting
a unique predicate per statement. Both are genuine improvements over
classical reification for the \emph{annotation} use case, and both are
a different bet than ours. RDF-star keeps the triple as the storage
unit and adds a nesting mechanism on top: a fact of arity $k$ is still
$k$ triples plus quoted-triple annotations, an entity-anchored query
still reassembles the fact through joins, and context remains an
annotation on statements rather than an index over them. SFDB inverts
that bet: the atomic storage unit is the whole $n$-ary fact, and the
context poset is the primary index, so entity-anchored retrieval is a
stalk lookup rather than a join, at the memory cost quantified in
Section~\ref{sec:eval:memory}. Whether RDF-star's engine-level
implementations (which are maturing quickly) close the retrieval gap in
practice is an empirical question our evaluation does not answer; the
architectural distinction --- annotate statements versus index by
context --- is the point this paper explores.

\paragraph{Labeled property graphs.}
\label{sec:related-work:lpg}
Labeled property graphs (LPGs)~\cite{angles2017property}, as
implemented by Neo4j, Amazon Neptune, and TigerGraph, are the other
obvious answer to ``why not just use a system that already handles
$n$-ary facts.'' An LPG attaches a key-value property map directly to
each node and edge, so a binary relation with metadata (an
\texttt{employed\_at} edge carrying a \texttt{since} property) needs no
reification at all. This genuinely dissolves the reification problem
for facts whose natural arity is two: subject and object are the
edge's endpoints, and everything else is a property on that one edge,
retrieved with the edge in a single lookup. It does not dissolve it for
facts of higher arity, which is the case this paper is about: a
property map is attached to one edge between two nodes, so a fact with
three or more essential participants (a grant awarded to a recipient by
a funder for a project) still needs either a synthetic intermediate
node with several edges radiating from it --- the same reification this
paper's KG baseline performs, just drawn as a graph instead of a triple
store --- or an awkward encoding as edge properties on one arbitrarily
chosen pair of participants, which asymmetrically privileges two of the
$n$ participants over the rest. Context is a second, independent gap:
LPG query languages (Cypher, Gremlin) have no native notion of a
context poset to index by. Scoping in practice is done either with
property filters evaluated per node/edge at query time (linear in the
matched set, not a direct lookup) or with separate physical databases
per tenant, neither of which offers the direct \texttt{context:<c>}
open-set lookup or the restriction-map structure this paper's context
poset provides (Section~\ref{sec:formal_definitions}). The comparison
is therefore not ``sheaf-indexed store versus property graph'' in the
abstract; it is that LPGs solve a real subset of the reification
problem this paper opens with --- binary facts with metadata --- cleanly
and natively, while leaving both higher-arity reification and
hierarchical context indexing as open problems on top of the same
substrate this paper's KG baseline already represents. We do not
benchmark against Neo4j or Neptune directly: our KG baseline's
reified-triple SPO/POS/OSP design is architecturally closer to the
production RDF stores we do compare against
(Section~\ref{sec:eval:jena}--Section~\ref{sec:eval:graphdb}) than to an
LPG's native edge-property storage, and a fair LPG comparison would
require a fourth adapter and a schema translation this paper does not
attempt; we name it here as the most natural fourth-store extension
(Section~\ref{sec:future:jena}) rather than claim a result we did not
measure.

\paragraph{Categorical and topological database models.}
The categorical database literature~\cite{spivak2012functorial} recasts
schemas as categories and instances as functors. Sheaf structures have
been used for distributed sensor fusion~\cite{robinson2017sheaves} and
for reasoning about locally consistent but globally inconsistent data
collections, where obstructions to global sections characterise
contextuality~\cite{abramsky2011sheaf}. Simplicial
complexes and cellular sheaves have been proposed as data models for
higher-order relationships~\cite{edelsbrunner2010computational}. Formal Concept
Analysis~\cite{ganter1999formal} organises data around a lattice of concepts
that shares surface structural features with our context poset, but does not
include the restriction maps or stalks that our design uses operationally.
Incidence algebras over locally finite posets~\cite{rota1964moebius} are,
in our own survey of candidate formalisms, the nearest mathematical
neighbour to the context-poset design: they support M{\"o}bius-inversion
style aggregation over exactly the poset structure we index by, and we
regard an incidence-algebra query layer as the most promising follow-up
formalism (Section~\ref{sec:future}).

We borrow terminology and structural intuition from this literature, but
we are not proving a theorem about sheaves. As we make explicit throughout, the
sheaf condition is vacuous on the finite Alexandrov topology that fits our
use case, so anything the model buys us empirically is buying it because of
the presheaf structure and the choice of indexing, not because of the sheaf
axioms.

\paragraph{Context and provenance in semantic data.}
Contexts, named graphs, and provenance annotations have received sustained
attention in the semantic web community~\cite{carroll2005named}. Approaches
range from explicit named graphs, to context-carrying quads, to reified
provenance ontologies such as PROV-O. Our contribution is neither a new
context ontology nor a new provenance model; we adopt a straightforward
dot-separated hierarchy of contexts as our first-class organising principle and
show what a database built around it looks like when the mathematical structure
is taken seriously.

\paragraph{Positioning.}
To the best of our knowledge, no prior system provides all of the following
together: a context-indexed semantic fact storage engine with stalk-based
LOOKUP, a co-implemented reified triple-store baseline, a canonical
intermediate representation with a constructive semantic-equivalence proof
between the two that a runtime harness checks on every insert and query, and
a benchmark evaluation that reports both wins and losses
across a LOOKUP/GLOBAL/CONTEXT/TEMPORAL taxonomy. The individual ingredients are all standard,
in some cases decades old. What this paper contributes is a coherent design
that combines them and an honest evaluation of what the combination is and is
not worth.
